% ===================================================================
%  Dodatek E: Kwantyzacja pola Φ w~duchu TGP
%  Schemat samozwrotny — bez tła przestrzennego
% ===================================================================
\section{Kwantyzacja pola $\Phi$: schemat samozwrotny}%
\label{app:kwantyzacja}

\subsection{Problem: dlaczego standardowa QFT nie stosuje się}%
\label{app:E-problem}

Standardowa kwantowa teoria pola (QFT) kwantuje pola
\textbf{na danym tle przestrzennym}: wybieramy rozmaitość
$(M, g_{\mu\nu})$, rozkładamy pole na mody normalne,
definiujemy próżnię Focka, itp. Podejście to zakłada,
że przestrzeń istnieje jako \emph{uprzednio dana arena}.

W~TGP ta droga jest zamknięta z~przyczyn ontologicznych
(aksjomat~\ref{ax:przestrzen}): $\Phi$ konstytuuje przestrze\'n,
nie ,,mieszka'' na niej. Nie możemy zdefiniować próżni Focka
dla $\Phi$ na tle metryki, która sama zależy od $\Phi$.

\begin{remark}[Pętla samozwrotna]
W~TGP zachodzi pętla:
\[
    \Phi \;\longrightarrow\; g_{\mu\nu}(\Phi)
    \;\longrightarrow\; \text{kwantyzacja na tle }g
    \;\longrightarrow\; \langle\delta\Phi^2\rangle
    \;\longrightarrow\; \Phi.
\]
Pętla ta nie jest defektem --- jest \emph{istotą} teorii.
Kwantyzacja musi być \textbf{samozwrotna}: fluktuacje $\Phi$
muszą być zdefiniowane względem struktury generowanej przez
samego $\Phi$.
\end{remark}

% -------------------------------------------------------------------
\subsection{Schemat kwantyzacji substratowej}%
\label{app:E-scheme}
% -------------------------------------------------------------------

Zamiast kwantyzować pole $\Phi$ na przestrzeni,
kwantyzujemy \textbf{substrat} $\Gamma = (V, E)$ ---
obiekt fundamentalny, niezależny od metryki.

\begin{definition}[Przestrzeń Hilberta substratu]%
\label{def:hilbert-substrate}
Przestrzenią Hilberta substratu jest:
\begin{equation}\label{eq:H-substrate}
    \mathcal{H}_\Gamma
    = \bigotimes_{i \in V} \mathcal{H}_i,
    \qquad
    \mathcal{H}_i = L^2([-1, 1], d\sigma_i),
\end{equation}
gdzie $d\sigma_i$ jest miarą Lebesgue'a na stanie węzła $i$.
Obserwable są operatorami samozwrotnymi na $\mathcal{H}_\Gamma$.
\end{definition}

Pole $\Phi$ w~tym schemacie jest \textbf{obserwalnym blokowym}:
\begin{equation}\label{eq:Phi-observable}
    \hat{\Phi}(x) = \frac{1}{N_B}
    \sum_{i \in \mathcal{B}(x)} \hat{s}_i^2
    \;\in\; \mathcal{B}(\mathcal{H}_\Gamma).
\end{equation}

\begin{proposition}[Próżnia substratowa]\label{prop:substrate-vacuum}
Próżnia TGP jest \textbf{stanem termalnym substratu}
w~fazie metrycznej ($T < T_c$):
\begin{equation}\label{eq:vacuum-thermal}
    |0_{\mathrm{TGP}}\rangle
    \equiv \rho_0
    = \frac{e^{-H_\Gamma/k_B T_{\mathrm{sub}}}}
           {\mathrm{Tr}\,e^{-H_\Gamma/k_B T_{\mathrm{sub}}}},
    \qquad T_{\mathrm{sub}} < T_c.
\end{equation}
Wartość oczekiwana:
$\langle\hat{\Phi}(x)\rangle_{\rho_0} = \PhiZero = v^2$
(par.\ porządku), co odpowiada jednorodnej próżni TGP.
\end{proposition}

% -------------------------------------------------------------------
\subsection{Propagator kwantowy pola $\Phi$}%
\label{app:E-propagator}
% -------------------------------------------------------------------

\begin{definition}[Fluktuacja substratowa]\label{def:Phi-fluct}
Fluktuacją kwantową pola $\Phi$ nazywamy:
\begin{equation}\label{eq:delta-Phi}
    \delta\hat{\Phi}(x)
    = \hat{\Phi}(x) - \PhiZero,
\end{equation}
z~operatorem $\hat{\Phi}(x)$ określonym przez~\eqref{eq:Phi-observable}.
\end{definition}

\begin{theorem}[Propagator kwantowy TGP]\label{thm:propagator}
W~granicy ciągłej i~w~przybliżeniu gaussowskim (drugie pochodne
działania) propagator Feynmana dla $\delta\Phi$ wynosi:
\begin{equation}\label{eq:propagator}
    G_\Phi(x, y)
    = \langle T[\delta\hat{\Phi}(x)\,\delta\hat{\Phi}(y)]\rangle_{\rho_0}
    = \int \frac{d^4k}{(2\pi)^4}\;
      \frac{i}{k^2 - m_{\mathrm{sp}}^2 + i\epsilon}\;
      e^{ik(x-y)},
\end{equation}
gdzie $m_{\mathrm{sp}}^2 = \gamma > 0$ jest masą efektywną
z~stw.~\ref{prop:vacuum-stability}.
\end{theorem}

\begin{proof}
Gaussowska aproksymacja działania wokół $\Phi = \PhiZero$
(zlinearyzowane równanie pola, stw.~\ref{prop:vacuum-stability})
daje działanie Kleina--Gordona:
\begin{equation}
    S_{\mathrm{KG}}[\delta\Phi]
    = \frac{1}{2}\int\!\left[
        (\partial_\mu\delta\Phi)^2
        - m_{\mathrm{sp}}^2\,(\delta\Phi)^2
    \right]\sqrt{-g}\,d^4x,
\end{equation}
z~$g_{\mu\nu} = \eta_{\mu\nu}$ w~przybliżeniu płaskiej próżni.
Standardowy propagator Klein--Gordona daje~\eqref{eq:propagator}.
Masa $m_{\mathrm{sp}}^2 = \gamma > 0$ zapewnia biegun rzeczywisty
(masywne kwanty przestrzeni).
\end{proof}

\begin{remark}[Propagator samozwrotny]\label{rem:self-ref-prop}
Propagator~\eqref{eq:propagator} jest zdefiniowany na \emph{płaskim}
tle $\eta_{\mu\nu}$ --- to jest \textbf{przybliżenie}: zakłada,
że fluktuacje $\delta\Phi$ są małe i~nie zmieniają metryki.
Pełna teoria wymaga iteracyjnego schematu:
\begin{enumerate}[nosep]
    \item Startujemy z~$g_{\mu\nu}^{(0)} = \eta_{\mu\nu}$.
    \item Obliczamy propagator $G_\Phi^{(0)}(x,y)$.
    \item Obliczamy wkład pętlowy do $\langle\Phi\rangle$:
          $\delta\Phi^{(1)} = \hbar \cdot G_\Phi^{(0)}(x,x)$.
    \item Aktualizujemy metrykę: $g^{(1)}_{\mu\nu}[\Phi^{(0)}
          + \delta\Phi^{(1)}]$.
    \item Powtarzamy aż do zbieżności.
\end{enumerate}
Na poziomie zera pętlowego ($\hbar \to 0$) odzyskujemy klasyczne
równania TGP. Na poziomie jednej pętli pojawiają się kwantowe
korekcje do $\Phi$, $\Lambda_{\mathrm{eff}}$ i~stałych sprzężenia.
\end{remark}

\begin{proposition}[Pierwsza iteracja schematu samozwrotnego]%
\label{prop:first-iteration}
Jawny wynik \textbf{pierwszej iteracji} (n=1) schematu
z~uwagi~\ref{rem:self-ref-prop}:

\medskip
\noindent\textbf{(1) Propagator zerowego rzędu.}
Na tle $g_{\mu\nu}^{(0)} = \eta_{\mu\nu}$, po obróceniu Wicka
($k^0 \to ik_E^0$) propagator euklidesowy:
\begin{equation}\label{eq:prop-eucl}
    G_\Phi^{(0)}(x,x)\big|_{\mathrm{E}}
    = \int_0^{\Lambda_{\mathrm{UV}}} \frac{d^4k_E}{(2\pi)^4}
      \frac{1}{k_E^2 + m_{\mathrm{sp}}^2}
    = \frac{1}{16\pi^2}
      \left[\Lambda_{\mathrm{UV}}^2 - m_{\mathrm{sp}}^2\,
      \ln\!\left(\frac{\Lambda_{\mathrm{UV}}^2}{m_{\mathrm{sp}}^2} + 1\right)\right],
\end{equation}
z~$\Lambda_{\mathrm{UV}} = \lP^{-1}$ i~$m_{\mathrm{sp}}^2 = \gamma$.
W~granicy $\Lambda_{\mathrm{UV}} \gg m_{\mathrm{sp}}$:
\begin{equation}\label{eq:prop-approx}
    G_\Phi^{(0)}(x,x)\big|_{\mathrm{reg}}
    \approx \frac{\lP^{-2}}{16\pi^2}
    - \frac{\gamma}{16\pi^2}\ln\!\frac{1}{\gamma\lP^2}.
\end{equation}

\noindent\textbf{(2) Kwantowa korekcja do próżni $\Phi$.}
Pierwszy iteracyjny wkład do średniej wartości pola:
\begin{equation}\label{eq:Phi-1st-corr}
    \delta\Phi^{(1)} = \hbar_0\,G_\Phi^{(0)}(x,x)\big|_{\mathrm{reg}}
    \approx \frac{\hbar_0\,\lP^{-2}}{16\pi^2}
    \left[1 - \gamma\lP^2\ln\!\frac{1}{\gamma\lP^2}\right].
\end{equation}
Wyrażenie $\gamma\lP^2 \sim 10^{-52}\times 10^{-70} = 10^{-122} \ll 1$,
więc człon logarytmiczny jest pomijalny, a~$\delta\Phi^{(1)}
\approx \hbar_0\lP^{-2}/(16\pi^2)$.

\noindent\textbf{(3) Zaktualizowana metryka.}
Z~hipotezy~\ref{hyp:metric-exp}: $g_{\mu\nu}^{(1)} =
g_{\mu\nu}[\PhiZero + \delta\Phi^{(1)}]$.
Perturbacja $\delta\Phi^{(1)}/\PhiZero \sim \hbar_0 \lP^{-2}/(16\pi^2\PhiZero)
\sim \lP^{-2}/\PhiZero \cdot \hbar_0$.
W~jednostkach Plancka ($\hbar_0 = 1$, $\lP = 1$):
$\delta\Phi^{(1)}/\PhiZero \sim 1/(16\pi^2\PhiZero) \sim 10^{-3}$.
Korekcja metryczna jest zatem supresowana przez $1/(16\pi^2\PhiZero)$.

\noindent\textbf{(4) Korekcja do stałej kosmologicznej.}
Z~propozycji~\ref{prop:loop-Lambda}: korekcja pętlowa jest
$\delta\Lambda^{(1)} = \frac{m_{\mathrm{sp}}^2}{2}
G_\Phi^{(0)}(x,x)|_{\mathrm{reg}} \approx
\frac{\gamma}{32\pi^2}\left[\lP^{-2} - \gamma\ln({\lP^{-2}/\gamma})\right]$.
Dominuje człon $\frac{\gamma\lP^{-2}}{32\pi^2}$, ale:
\begin{equation}\label{eq:dLambda-ratio}
    \frac{\delta\Lambda^{(1)}}{\Lambda_{\mathrm{eff}}}
    = \frac{\gamma\lP^{-2}/(32\pi^2)}{\gamma/12}
    = \frac{12}{32\pi^2\lP^2\cdot 1}
    \approx \frac{12}{32\pi^2} \cdot \frac{1}{\lP^2\cdot\gamma^{-1}}\,.
\end{equation}
Ponieważ $\lP^2 \cdot \gamma = \lP^2 \cdot \Lambda_{\mathrm{obs}} \sim 10^{-122}$,
stąd $\delta\Lambda^{(1)}/\Lambda_{\mathrm{eff}} \sim 12\lP^{-2}/(32\pi^2) \cdot \lP^2/\Lambda_{\mathrm{obs}}^{-1}$
--- i~powracamy do problemu hierarchii, \textbf{chyba że} przyjmiemy,
że $\Lambda_{\mathrm{UV}}^{\mathrm{eff}} = \sqrt{\gamma}$ (cięcie IR,
nie UV), co jest spójne z~masą pola:
\begin{equation}\label{eq:IR-cutoff-argue}
    \Lambda_{\mathrm{UV}}^{\mathrm{eff}} \to m_{\mathrm{sp}} = \sqrt{\gamma},
    \qquad
    \delta\Lambda^{(1)}_{\mathrm{IR}}
    = \frac{\gamma}{32\pi^2}\cdot\gamma\cdot
    [\text{logarytm}] \propto \gamma^2.
\end{equation}
Ta subtelność --- wybór skali regulatora --- jest
\textbf{otwartym problemem~O15}: pełna iteracyjna zbieżność
schematu samozwrotnego wymaga rachunku z~nieperturbatywnym
obcięciem na skali substratowej.
\end{proposition}

% -------------------------------------------------------------------
\subsection{Regularyzacja: naturalna obcięcie w~$\lP$}%
\label{app:E-cutoff}
% -------------------------------------------------------------------

\begin{theorem}[Naturalny cutoff kwantowy]\label{thm:natural-cutoff}
W~TGP nie potrzeba zewnętrznego parametru regularyzacyjnego.
Naturalnym cutoffem UV jest \textbf{skala Plancka}
$\Lambda_{\mathrm{UV}} = \lP^{-1}$, wynikająca ze skali kratowej
substratu:
\begin{equation}\label{eq:UV-cutoff}
    a_{\mathrm{sub}} = \lP
    \;\;\Longrightarrow\;\;
    \Lambda_{\mathrm{UV}}^{-1} = \lP = \sqrt{\frac{\hbar_0 G_0}{c_0^3}}.
\end{equation}
Twierdzenie~\ref{thm:lP} gwarantuje, że $\lP$ jest stałą ---
cutoff jest \textbf{absolutny}, nie zależy od $\Phi$.
\end{theorem}

\begin{proof}
Substrat $\Gamma$ ma skala kratową $a_{\mathrm{sub}}$.
Coarse-graining ($a_{\mathrm{sub}} \to 0$ w~granicy ciągłej)
jest cutoffem dla teorii efektywnej $\Phi$.
Identyfikacja $a_{\mathrm{sub}} = \lP$ wymaga trzech kroków:

\textbf{(i) Dlaczego jakaś stała skala?}
Twierdzenie~\ref{thm:lP}: $\lP = \sqrt{\hbar_0 G_0/c_0^3}$ jest
\textbf{stałą} w~TGP, ponieważ iloraz $\hbar_0 G_0/c_0^3$ nie zależy
od $\Phi$ (stałe $\hbar_0$, $G_0$, $c_0$ są wartościami referencyjnymi,
a~$\lP$ jest przez nie skonstruowana jako kombinacja niezmiennicza
skalowania $\Phi$-zależnego, por.\ sekcja~\ref{sec:stale}).
Substrat musi mieć \emph{jakąś} stałą skalę kratową --- fizyka
mikro-substratu jest odcięta od dynamiki makroskopowej $\Phi$.

\textbf{(ii) Dlaczego $a_{\mathrm{sub}} = \lP$, a~nie inna wartość?}
W~grawitacji kwantowej (pętlowej, strunowej, spinowych pianach)
skala dyskretu przestrzennego jest niezmienniczo związana z~$\lP$
przez warunek, że elementy objętości substratu mają minimalną
nieredukowalną objętość $\sim \lP^3$.
W~TGP: $\Phi > 0$ oznacza istnienie przestrzeni; minimalne
nieredukowalne $\delta\Phi = \PhiZero \cdot (\lP/a_{\mathrm{sub}})^3$
przy $\lP = a_{\mathrm{sub}}$ daje $\delta\Phi \sim \PhiZero$ ---
kwant przestrzenności rzędu jeden (bezwymiarowy), co jest
spójne z~próżniowym $\PhiZero \approx 25$.

\textbf{(iii) Brak dodatkowego skali.}
Inne stałe długości w~TGP (np.\ $m_{\mathrm{sp}}^{-1} = \gamma^{-1/2}
\sim 3\;\mathrm{Mpc}$) są dynamiczne i~kosmologiczne --- nie mogą
być skalą kratową mikro-fizyki. Jedyną pozostałą skalą jest~$\lP$.

Wniosek: $a_{\mathrm{sub}} = \lP$, co daje~\eqref{eq:UV-cutoff}.
\end{proof}

\begin{corollary}[Skończona renormalizacja]\label{cor:finite-renorm}
Przy cutoffie $\Lambda_{\mathrm{UV}} = \lP^{-1}$:
\begin{enumerate}[nosep]
    \item Dywergencje pętlowe są \textbf{mocami} $\lP^{-n}$,
          a~nie $\infty$ --- wszystkie całki są skończone.
    \item Renormalizowane parametry (masy, stałe sprzężenia)
          absorbują poprawki pętlowe w~skalę $\lP$ --- zgodnie
          z~standardową renormalizacją efektywną.
    \item Nie ma problemu hierarchii w~tradycyjnym sensie:
          $m_H^2 / m_{\mathrm{Planck}}^2 \sim v_W^2/\lP^{-2}$
          jest pytaniem o~substratowe parametry $\lambda_0, v$,
          nie o~fine-tuning.
\end{enumerate}
\end{corollary}

% -------------------------------------------------------------------
\subsection{Korekcje kwantowe do stałej kosmologicznej}%
\label{app:E-cosmconst}
% -------------------------------------------------------------------

\begin{proposition}[Wkład pętlowy do $\Lambda_{\mathrm{eff}}$]%
\label{prop:loop-Lambda}
Jednoopętlowy wkład kwantów $\Phi$ do efektywnej stałej
kosmologicznej:
\begin{equation}\label{eq:loop-Lambda}
    \delta\Lambda^{(1)}
    = \frac{1}{2}\,m_{\mathrm{sp}}^2\,
      G_\Phi(x,x)\big|_{\mathrm{reg}}
    = \frac{m_{\mathrm{sp}}^2}{8\pi^2}
      \left[\Lambda_{\mathrm{UV}}^2 - m_{\mathrm{sp}}^2
      \ln\frac{\Lambda_{\mathrm{UV}}^2}{m_{\mathrm{sp}}^2}\right],
\end{equation}
gdzie $\Lambda_{\mathrm{UV}} = \lP^{-1}$ i~$m_{\mathrm{sp}}^2 = \gamma$.
\end{proposition}

\begin{remark}[Naturalność $\Lambda_{\mathrm{eff}}$ w~TGP]%
\label{rem:naturalness}
Podstawiając $m_{\mathrm{sp}}^2 = \gamma \sim H_0^2/c_0^2
\sim 10^{-52}\;\mathrm{m}^{-2}$
i~$\Lambda_{\mathrm{UV}}^2 = \lP^{-2} \sim 10^{70}\;\mathrm{m}^{-2}$:
\begin{equation}
    \delta\Lambda^{(1)}
    \sim \frac{\gamma}{\lP^2} \cdot \frac{\lP^2}{8\pi^2}
    = \frac{\gamma}{8\pi^2}
    \approx \frac{10^{-52}}{80} \;\mathrm{m}^{-2},
\end{equation}
co jest rzędu obserwowanego $\Lambda_{\mathrm{obs}}$!

\textbf{Kluczowe}: w~TGP nie ma problemu hierarchii
$10^{122}$ między $\Lambda_{\mathrm{QFT}}$ a~$\Lambda_{\mathrm{obs}}$
dlatego, że:
\begin{itemize}[nosep]
    \item Masa pola $\Phi$ to $m_{\mathrm{sp}}^2 = \gamma$ ---
          ta sama stała, która wyznacza $\Lambda_{\mathrm{eff}} = \gamma/12$.
    \item Wkład pętlowy~\eqref{eq:loop-Lambda} jest
          \textbf{proporcjonalny do} $\gamma$, nie do $\lP^{-2}$.
    \item Kwanty $\Phi$ (fonony przestrzenności) mają masę
          wyznaczoną przez $\Lambda_{\mathrm{eff}}$, nie przez skalę
          Plancka --- to jest samospójna pętla, nie fine-tuning.
\end{itemize}
\end{remark}

% -------------------------------------------------------------------
\subsection{Stany wzbudzone: fonony przestrzenności}%
\label{app:E-phonons}
% -------------------------------------------------------------------

\begin{definition}[Fonon przestrzenności]\label{def:phonon}
\textbf{Fononem przestrzenności} nazywamy kwant fluktuacji
pola $\Phi$ wokół próżni $\PhiZero$:
\begin{equation}\label{eq:phonon-state}
    |k\rangle_\Phi = \hat{a}^\dagger(k)|0_{\mathrm{TGP}}\rangle,
    \qquad
    \hat{a}^\dagger(k) = \int \frac{d^3x}{(2\pi)^3}
    e^{-ikx}\,\hat{a}^\dagger_{\mathbf{k}},
\end{equation}
z~relacją dyspersji:
\begin{equation}\label{eq:phonon-dispersion}
    \omega_k^2 = c_0^2\,\mathbf{k}^2 + m_{\mathrm{sp}}^2\,c_0^4,
    \qquad
    m_{\mathrm{sp}} = \sqrt{\gamma}\,\frac{\hbar_0 c_0}{l_{\mathrm{sp}}}
    \;\sim\; \frac{\hbar_0 H_0}{c_0},
\end{equation}
gdzie $l_{\mathrm{sp}} = 1/\sqrt{\gamma} \approx R_H$
(zasięg Yukawy~\ref{eq:lambda-Yukawa}).
\end{definition}

\begin{remark}[Ultralekie fonony]
Masa fononu przestrzenności:
$m_{\mathrm{sp}}\,c^2 \sim \hbar\,H_0 \sim 10^{-33}\;\mathrm{eV}$
--- ekstremalne ultralekkie. Jest to naturalny kandydat na
\textbf{emergentną ciemną energię}: zbiorowe wzbudzenia
pola przestrzenności mają masę wyznaczoną przez $H_0$,
co tłumaczy zbieżność skali ciemnej energii ze skalą Hubble'a
(problem coincidence, hyp.~\ref{hyp:coincidence}).
\end{remark}

% -------------------------------------------------------------------
\subsection{Diagram Feynmana: TGP jako QFT samozwrotna}%
\label{app:E-feynman}
% -------------------------------------------------------------------

\begin{remark}[Reguły Feynmana TGP]\label{rem:feynman-rules}
Diagramy Feynmana dla TGP mają niestandardową interpretację:
\begin{enumerate}[nosep]
    \item \textbf{Propagator fononu} $\Phi$: linia przerywana,
          wartość $G_\Phi(k) = i/(k^2 - \gamma + i\epsilon)$.
    \item \textbf{Vertex wierzchołkowy} $\Phi^3$ (kubiczny):
          wynika z~członu $\beta\Phi^3/\PhiZero^2$, siła
          sprzężenia $\sim \beta/\PhiZero$.
    \item \textbf{Vertex kwartyczny} $\Phi^4$:
          wynika z~$-\gamma\Phi^4/\PhiZero^3$,
          siła $\sim \gamma/\PhiZero^2$.
    \item \textbf{Sprzężenie z~materią}: vertex $\Phi$--materia
          wynika z~członu źródłowego $q\rho$ w~równaniu pola;
          amplituda $\sim q$.
    \item \textbf{Fotony / bozony cechowania}: linie ciągłe
          z~propagatorem standardowym, ale na tle efektywnym
          $g_{\mu\nu}(\Phi)$.
\end{enumerate}

\textbf{Niestandardowość}: w~standardowej QFT na tle grawitacyjnym
linie materii i~grawitony są rozdzielone. W~TGP \emph{nie ma}
,,grawitonu'': propagatorem grawitacji jest sam propagator $\Phi$.
Foton $A_\mu$ propaguje na tle $g_{\mu\nu}(\Phi)$, które zależy
od $\Phi$ --- dlatego diagramy $\Phi$-$A_\mu$ mogą generować
niekonwencjonalne korekcje do propagacji fotonu.
\end{remark}

% -------------------------------------------------------------------
\subsection{Problem otwarty: pełna kwantyzacja}%
\label{app:E-open}
% -------------------------------------------------------------------

\begin{problem}[Kwantyzacja pola $\Phi$ --- OTWARTY]%
\label{prob:kwantyzacja}
Niniejszy dodatek dostarcza:
\begin{enumerate}[nosep]
    \item \textbf{Schemat}: kwantyzacja substratu jako przestrzeni
          Hilberta $\mathcal{H}_\Gamma$.
    \item \textbf{Propagator} w~przybliżeniu gaussowskim.
    \item \textbf{Cutoff}: $\lP$ jako skala kratowa.
    \item \textbf{Korekcje pętlowe} do $\Lambda_{\mathrm{eff}}$.
\end{enumerate}
Pozostaje otwarte:
\begin{itemize}[nosep]
    \item Pełna teoria perturbacji (wyższe pętle,
          wkłady nieperturbacyjne instantonów substratowych).
    \item Samospójna iteracja $\Phi \to g_{\mu\nu} \to G_\Phi
          \to \delta\langle\Phi\rangle \to g_{\mu\nu}$
          (problem O15).
    \item Unitarność: czy $S$-macierz TGP jest unitarna
          mimo samozwrotności?
    \item Statystyki: fermiony w~TGP jako topologiczne defekty
          --- czy statystyki Fermiego--Diraca wynikają
          z~topologii profilu radialnego (prob.~\ref{prob:fermion-full})?
\end{itemize}
To jest \textbf{program otwartej kwantyzacji TGP} (O15).
\end{problem}

% -------------------------------------------------------------------
\subsection{Cz\k{a}stki wirtualne w~TGP (O11)}%
\label{app:E-virtual}
% -------------------------------------------------------------------

\begin{definition}[Cz\k{a}stka wirtualna w~TGP]\label{def:virtual-particle}
\textbf{Cz\k{a}stk\k{a} wirtualn\k{a}} nazywamy lokaln\k{a} fluktuacj\k{e}
pola $\Phi$ poni\.zej progu pr\'o\.zni, tj.\ zdarzenie
\begin{equation}\label{eq:virtual-condition}
    \Phi(\mathbf{x}, t) < \PhiZero
    \quad \Longleftrightarrow \quad
    \chi(\mathbf{x}, t) < 1,
\end{equation}
kt\'ore jest czasowo-przestrzennie ograniczone
(ograniczone do obszaru $|\Delta\mathbf{x}| \lesssim \Delta t \cdot c_0$)
i~nie rozprzestrzenia si\k{e} asymptotycznie jako stan rzeczywisty.

Geometrycznie: fluktuacja $\chi < 1$ odpowiada lokalnej
\emph{nicości} zanurzanej w~przestrzeni~---
fragmentowi fazy $\mathcal{S}_0$
zanurzanemu chwilowo w~substrat fazy $\mathcal{S}_1$.
\end{definition}

\begin{remark}[Interpretacja i~por\'ownanie ze standardow\k{a} QFT]
\label{rem:virtual-comparison}
W standardowej QFT cz\k{a}stki wirtualne s\k{a} wewn\k{e}trznymi
liniami diagram\'ow Feynmana --- formalnym narz\k{e}dziem
rachunkowym, bez interpretacji fizycznej.

W~TGP uzyskuj\k{a} \textbf{bezpo\'sredni\k{a} interpretacj\k{e} ontologiczn\k{a}}:
\begin{enumerate}[nosep]
    \item \textbf{Wirtualny foton} ($\Phi < \PhiZero$ w~obszarze
          mi\k{e}dzy \l{}adunkami): lokalna redukcja g\k{e}sto\'sci
          przestrzenno\'sci transmituje oddzia\l{}ywanie.
    \item \textbf{Masy wirtualne}: fluktuacja $\chi < 1$ powoduje
          $c(\Phi) > c_0$ (prawo ax.~\ref{ax:c}),
          $G(\Phi) < G_0$ (ax.~\ref{ax:G}) ---
          cz\k{a}stka wirtualna zmienia lokalnie prawa fizyki.
    \item \textbf{Czas \.zycia}: z~zasady nieoznaczono\'sci w~TGP
          $\Delta t \cdot \Delta E \sim \hbar_0$,
          przy czym $\hbar_0$ jest g\'orn\k{a} granic\k{a} produktu
          (emergentne $\hbar(\Phi) \leq \hbar_0$).
\end{enumerate}
\end{remark}

\begin{proposition}[Warunek istnienia cz\k{a}stki wirtualnej]\label{prop:virtual-existence}
Fluktuacja $\chi < 1$ jest dynamicznie dozwolona,
gdy energia fluktuacji spe\l{}nia:
\begin{equation}\label{eq:virtual-energy}
    \delta E \;=\; \int_{\delta V}
    \Bigl[U(\chi) - U(1)\Bigr]\,d^3x
    \;\approx\; \PhiZero^2\!\int_{\delta V}
    \frac{(\chi-1)^2}{2}\Bigl(\gamma - \beta\Bigr)\,d^3x
    \;\leq\; \frac{\hbar_0}{\Delta t},
\end{equation}
gdzie $\delta V$ to obj\k{e}to\'s\'c fluktuacji,
a~warunek $\gamma > \beta$ zapewnia,
\.ze pr\'o\.znia $\chi=1$ jest \emph{lokalnym minimum} potencjalu
(co jest spe\l{}nione przy naturalnych WP:
$\gamma = \beta$ daje stan brzegin; $\gamma > \beta$ --- stabil.).
\end{proposition}

\begin{remark}[Status O11]
Powy\.zsza definicja i~propozycja stanowi\k{a} \textbf{cz\k{e}\'sciowe zamkni\k{e}cie}
problemu O11.
Otwarte pozostaje: pe\l{}na rachunkowp\'os\'c diagram\'ow Feynmana
z~wirtualnymi cz\k{a}stkami, w~szczeg\'olno\'sci obliczenie
korekcji do oddzia\l{}ywania Coulomba i~Youkawy z~u\.zyciem
regu\l{} sekcji~\ref{app:E-feynman}
w~bazie ontologicznej TGP.
\end{remark}

% -------------------------------------------------------------------
\subsection{Spin-$\tfrac{1}{2}$ z~topologii kinku --- szkic (II.A)}%
\label{app:E-spin-half}
% Sesja v32, 2026-03-24 --- zadanie II.A planu \texttt{PLAN\_ROZWOJU\_v1}
% -------------------------------------------------------------------

Niniejsza sekcja zawiera \textbf{szkic formalizmu} prowadz\k{a}cy
do spin-$\tfrac{1}{2}$ i~statystyk Fermiego--Diraca z~topologii
solitonu TGP.
Jest to odpowied\'z na problem~\ref{prob:fermion-full}
(por.\ prob.~\ref{prob:kwantyzacja}, sekcja~\ref{app:E-open}).

\begin{definition}[Przestrze\'n konfiguracyjna solitonu TGP]%
\label{def:kink-config-space}
\textbf{Przestrzeni\k{a} konfiguracyjn\k{a}} solitonu TGP nazywamy:
\begin{equation}\label{eq:kink-config}
  \mathcal{C}_{\mathrm{sol}}
  \;:=\;
  \bigl\{
    \varphi : \mathbb{R}^3 \to [0,\infty)
    \;\big|\;
    \varphi(0) = 0,\;
    \varphi(\infty) = 1,\;
    \text{sferycznie symetryczna}
  \bigr\} \!\big/ \!\sim,
\end{equation}
gdzie $\sim$ oznacza r\'ownowa\.zno\'s\'c przez ci\k{a}g\l{}e deformacje
zachowuj\k{a}ce warunki brzegowe.
Rozr\'o\.zniamy dwie klasy (por.\ def.~\ref{prop:kink-chirality}):
\begin{itemize}[nosep]
  \item \textbf{Kink}: $\varphi$ ro\'snie monotonicznie od $0$ do $1$;
  \item \textbf{Antykink}: transformacja $\varphi \to 2 - \varphi$
    (zamiana centrum i~pr\'o\.zni).
\end{itemize}
\end{definition}

\begin{proposition}[Indeks topologiczny kinku i~spin-$\tfrac{1}{2}$
  --- szkic]%
\label{prop:kink-spin-half}
Nast\k{e}puj\k{a}cy ci\k{a}g argument\'ow prowadzi do wniosku,
\.ze kink TGP niesie \textbf{spin-$\tfrac{1}{2}$}:
\begin{enumerate}[(S1)]
  \item \textbf{$\mathbb{Z}_2$-symetria substratu.}
    Hamiltoniano substratu~\eqref{eq:B-H} jest niezmienniczy
    wzgl\k{e}dem globalnej $\hat{s}_i \to -\hat{s}_i$.
    W~fazie uporzadkowanej $\langle\hat{s}^2\rangle \neq 0$:
    symetria spontanicznie z\l{}amana, lecz \emph{nieobserwowalna}
    dla $\Phi = \langle\hat{s}^2\rangle$ ($\mathbb{Z}_2$-parzyste).

  \item \textbf{Reprezentacja spinorowa.}
    Zmienne substratowe $\hat{s}_i$ transformuj\k{a} si\k{e}
    w~reprezentacji podw\'ojnej $SU(2) \twoheadrightarrow SO(3)$.
    Obrót o~$2\pi$ daje
    \begin{equation}\label{eq:Z2-rotation}
      R_{2\pi}|\Psi_{\mathrm{sub}}\rangle
      \;=\; (-1)^{N_s}\,|\Psi_{\mathrm{sub}}\rangle,
    \end{equation}
    gdzie $N_s$ jest liczb\k{a} w\k{e}z\l{}\'ow solitonu.

  \item \textbf{Homotopia przestrzeni konfiguracyjnej.}
    Po kompaktyfikacji ($\mathbb{R}^3 \cup \{\infty\} \cong S^3$)
    i~redukcji do symetrii sferycznych:
    \begin{equation}\label{eq:pi1-config}
      \pi_1(\mathcal{C}_{\mathrm{sol}}) \;=\; \mathbb{Z}_2.
    \end{equation}
    Nietrywialny element $\mathbb{Z}_2$ to p\k{e}tla obrotu
    o~$2\pi$ niecofalaj\k{a}ca si\k{e} do punktu
    (analogia Diraca: ska\'nca\l{}a wolna lina).

  \item \textbf{Znak stanu kinku.}
    Z~(S1)--(S3): przy obrocie o~$2\pi$ wok\'o\l{} dowolnej osi
    \begin{equation}\label{eq:kink-sign}
      \boxed{
        R_{2\pi}|\mathrm{kink}\rangle \;=\; -|\mathrm{kink}\rangle.
      }
    \end{equation}
    To jest w\l{}asno\'s\'c definiuj\k{a}ca \textbf{spin-$\tfrac{1}{2}$}.

  \item \textbf{Statystyki Fermiego--Diraca.}
    Wymiana dw\'och identycznych kink\'ow
    = obrót jednego o~$2\pi$ wzgl\k{e}dem drugiego.
    Z~\eqref{eq:kink-sign}:
    \begin{equation}\label{eq:exchange-sign}
      |\mathrm{kink}_1,\mathrm{kink}_2\rangle
      \;\xrightarrow{\text{wymiana}}\;
      -|\mathrm{kink}_2,\mathrm{kink}_1\rangle,
    \end{equation}
    sk\k{a}d \textbf{zasada Pauliego} i~statystyki FD.
\end{enumerate}
\end{proposition}

\begin{proposition}[Chiralno\'s\'c z~asymetrii kink--antykink]%
\label{prop:kink-chirality}
Kink i~antykink TGP tworz\k{a} dwie nierównoważne klasy
chiralne:
\begin{enumerate}[(C1)]
  \item \textbf{Kink-L} ($\varphi{:}\;0\!\to\!1$, profil narastaj\k{a}cy):
    odpowiada lewoskr\k{e}tno\'sci;
  \item \textbf{Kink-R} (antykink, $\tilde\varphi = 2-\varphi$):
    profil malej\k{a}cy od centrum,
    odpowiada prawoskr\k{e}tno\'sci.
\end{enumerate}
Asymetria pochodzi z~braku symetrii $\varphi \leftrightarrow 2-\varphi$
w~potencjale:
\begin{equation}\label{eq:chirality-asymmetry}
  U(\varphi) = \tfrac{\beta}{3}\varphi^3 - \tfrac{\gamma}{4}\varphi^4
  \;\;\not\equiv\;\;
  U(2-\varphi)
  \quad (\varphi \neq 1),
\end{equation}
co prowadzi do \textbf{chiralnego podzia\l{}u mas}
$m_L \neq m_R$ --- zal\k{a}\.zek mechanizmu Higgsa w~TGP.
\end{proposition}

\begin{problem}[Fermiony w~TGP: program pe\l{}nego wyprowadzenia]%
\label{prob:fermion-full}
Szkic tw.~\ref{prop:kink-spin-half}--\ref{prop:kink-chirality}
identyfikuje mechanizm.
Do pe\l{}nego zamkni\k{e}cia potrzebne s\k{a}:
\begin{enumerate}[nosep]
  \item Jawna kwantyzacja funkcjona\l{}u stanu solitonu
    w~$\mathcal{H}_\Gamma$ (def.~\ref{def:hilbert-substrate});
  \item Formalny dow\'od $\pi_1(\mathcal{C}_{\mathrm{sol}}) = \mathbb{Z}_2$
    dla sferycznie symetrycznych kink\'ow 3D;
  \item Twierdzenie spin--statystyki ze struktury
    algebry substratowej (\S\ref{app:E-scheme});
  \item Numeryczne obliczenie $m_L/m_R$ z~$U(\varphi)$.
\end{enumerate}
\textit{Status v32}: szkic jako\'sciowy zamkniety;
pe\l{}ne wyprowadzenie --- program Fazy~II.A
(\texttt{PLAN\_ROZWOJU\_v1}).
\end{problem}

\begin{remark}[Trzy generacje a~poziomy WKB]%
\label{rem:generations-wkb}
Trzy pokolenia fermion\'ow w~TGP odpowiadaj\k{a}
poziomom WKB solitonu: $n = 0, 1, 2$.
Ka\.zdy poziom niesie spin-$\tfrac{1}{2}$
z~tw.~\ref{prop:kink-spin-half} (niezale\.znie od~$n$).
Masy generacji wyznaczane s\k{a} przez warunek WKB
(\S\ref{hyp:generacje}), chiralny podzia\l{}
przez tw.~\ref{prop:kink-chirality}.
\textbf{Wniosek}: jednakowy spin wszystkich trzech generacji
jest konsekwencj\k{a} topologii, nie koincydencj\k{a}.
\end{remark}
